% Source: Wald GR, physical PDF p. 42 (printed p. 29).
% Coverage: Figure 3.1.
% Figures/tables/footnotes: Figure only; no tables or footnotes.
% Status: source-reviewed and compile-clean.
% Uncertainties: none.
\begingroup
\renewcommand{\thefigure}{3.1}
\begin{figure}[H]
  \centering
  \begin{tikzpicture}[scale=0.92,>=Stealth]
    \draw[thick] (-2.7,-0.8) -- (-1.0,0.8) -- (2.7,0.8) -- (1.2,-0.8) -- cycle;
    \draw[thick] plot[smooth cycle,tension=.75] coordinates
      {(-1.25,-0.1) (-.55,.40) (.82,.37) (1.25,-.08) (.56,-.37) (-.52,-.33)};
    \fill (-1.25,-0.1) circle (1.6pt) node[below left=1pt] {$p$};
    % The transported vector stays parallel to itself: every copy points
    % in the same (horizontal) direction, which is the point of the figure.
    \foreach \x/\y in {-.60/.42,.55/.40,1.28/-.06,.30/-.42,-.70/-.36} {
      \draw[->,thick] (\x,\y) -- ++(.42,0);
    }
    \draw[->,thick] (-1.20,-.10) -- ++(.42,0);
    \node[below=1pt] at (-.99,-.14) {$v^a$};
  \end{tikzpicture}
  \caption{The parallel transport of a vector, $v^a$, around a closed curve in the plane.
  The vector $v^a$ always ``comes back'' pointing in the same direction as it did initially.}
  \label{fig:3.1}
\end{figure}
\endgroup
